\documentclass[11pt,a4paper]{amsart}
\usepackage[T1]{fontenc}
\usepackage[margin=1in]{geometry}
\usepackage{amssymb}
\usepackage[hidelinks]{hyperref}
\hypersetup{pdftitle={Majority-SAT for 3-CNFs with One Unrestricted Clause}}
\DeclareMathSizes{8}{8}{8}{8}
\DeclareMathSizes{9}{9}{8}{8}
\DeclareMathSizes{10.95}{10.95}{8}{8}
\newtheorem{theorem}{Theorem}
\newtheorem{lemma}[theorem]{Lemma}
\newcommand{\PP}{\mathrm{PP}}
\newcommand{\sat}{\#\mathrm{SAT}}
\title{Majority-SAT for 3-CNFs with One Unrestricted Clause}
\date{November 27, 2026}

\begin{document}
\begin{abstract}
We prove PP-completeness, under polynomial-time many-one reductions, of testing whether a 3-CNF conjoined with a single unrestricted clause has satisfaction probability at least one half. The proof shifts an input integer independent-set threshold to a power of two. This also gives PP-completeness of comparing a 2-CNF's satisfaction probability with an input dyadic threshold of the form \(2^{-a}\). Akmal's selector transfers the latter result to majority, answering Open Problems 8 and 9 of his dissertation.
\end{abstract}
\maketitle

\section{The two comparisons}

Take Boolean assignments uniformly over the formula's variables. A \(k\)-CNF has width at most \(k\), and both comparisons below accept at equality. Akmal and Williams proved NP-hardness under Turing reductions for the one-long-clause majority problem \cite[Thm.~7.2]{akmalwilliams}. Akmal's dissertation asks for its exact complexity as Open Problem~8 and asks about the input power-of-two comparison on 2-CNFs as Open Problem~9 \cite[p.~94]{akmal}.

\begin{theorem}\label{main}
Each of these languages is \(\PP\)-complete under polynomial-time many-one reductions.
\begin{enumerate}
\item Given a 2-CNF \(\phi\) and a positive binary integer \(a\), decide if \(\Pr[\phi]\ge 2^{-a}\).
\item Given a 3-CNF \(\psi\) and one unrestricted clause \(L\), decide if \(\Pr[\psi\wedge L]\ge 1/2\).
\end{enumerate}
\end{theorem}

The result in part~1 uses an input exponent \(a\). Fixed probability thresholds for pure bounded-width CNFs are covered by Akmal and Williams \cite[Thm.~1.1]{akmalwilliams}. Finster et al.\ state PP-completeness of the independent-set threshold for general graphs \cite[App.~C.2, Problem~5, v4]{finster}. The weighted exclusion and dominant-count technique appears in Cozman and Mau\'a \cite[App.~A, proof of Thm.~2]{cozman}. Chan and Pak give the prescribed-count graph construction \cite[Ex.~5.8]{chanpak}. We combine a direct threshold proof with the graph join and Akmal's selector \cite[Lemma~4.11]{akmal}.

\section{From an integer threshold to majority}

Write \(I(B)\) for the number of independent sets of a simple graph \(B\), including the empty set. Take a graph \(B\) and a positive binary integer \(T\). Let \(b\) be the least nonnegative integer with \(2^b\ge T\) and set \(D=2^b-T+1\).

To construct the graph \(H\) of Chan and Pak \cite[Ex.~5.8]{chanpak} with \(I(H)=D\), start with the empty graph and read the binary digits of \(D\) after its leading \(1\). For each digit, add an isolated vertex to double the count. For a digit \(1\), then add a vertex adjacent to all existing vertices to increase the count by one. Thus \(H\) has \(O(\log D)\) vertices. Join \(B\) and \(H\) by adding every cross edge. Each independent set of the resulting graph \(J\) lies entirely in one side, with the empty set common to both. Therefore
\[
 I(J)=I(B)+D-1,\qquad I(J)\ge 2^b\ \Longleftrightarrow\ I(B)\ge T.
\]

Use a variable \(x_v\) for each vertex of \(J\) and impose \(\neg x_v\lor\neg x_w\) for each edge \(vw\). For every isolated \(v\), include
\((x_v\lor\neg w_v)\wedge(\neg x_v\lor\neg w_v)\) with a fresh \(w_v\). It mentions \(x_v\) explicitly and extends either value of \(x_v\) uniquely by \(w_v=0\). Append \(b+1\) fresh negative unit clauses on separate fresh variables. The resulting 2-CNF \(\phi\) has \(N\ge b+1\) occurring variables and exactly \(I(J)\) satisfying assignments. For \(a=N-b\ge1\),
\[
  \Pr[\phi]\ge 2^{-a}\ \Longleftrightarrow\ \sat(\phi)\ge 2^b
       \ \Longleftrightarrow\ I(B)\ge T .
\]
The number of new vertices, clauses and padding variables is polynomial in the graph size and the bit length of \(T\).

Apply Akmal's Lemma~4.11 \cite{akmal}. Add the same fresh literal \(s\) to each clause of \(\phi\) to form \(\psi\). Put \(L=(\neg s\lor y_1\lor\cdots\lor y_a)\) with fresh variables \(y_i\). Conditioning on \(s=0\) gives acceptance probability \(\Pr[\phi]\), and conditioning on \(s=1\) gives \(1-2^{-a}\). Thus
\[
 \Pr[\psi\wedge L]=\frac{\Pr[\phi]+1-2^{-a}}2 .
\]
In particular the maps preserve equality at the threshold.

\section{A direct proof of the source hardness}

\begin{lemma}\label{source}
Testing \(I(B)\ge T\) for a simple graph \(B\) and a positive binary integer \(T\) is \(\PP\)-hard under polynomial-time many-one reductions.
\end{lemma}

\begin{proof}
Start with a Boolean circuit \(C\) on \(r\ge1\) free bits and the test \(A\ge K\), where \(A\) is its accepting count and \(K=2^{r-1}\). This test is \(\PP\)-hard even if \(\PP\) is defined by strict majority. Given a circuit to be tested for strict majority, add one selector bit and use the input circuit in one branch and acceptance on exactly \(2^{r-1}-1\) strings in the other. The new count reaches half precisely when the old count strictly exceeded half. The auxiliary branch uses a binary comparator.

Encode the gates of \(C\) by their Boolean equivalences using clauses of width at most three and require the output to be true. For example, \(z=x\wedge y\) uses \((\neg z\lor x)\), \((\neg z\lor y)\), \((z\lor\neg x\lor\neg y)\). The dual rules encode OR and two clauses encode NOT. Each free input has a unique gate extension, so the resulting 3-CNF \(F\), with all free bits declared, has exactly \(A\) models.

Build a graph \(G\) with one group for each variable and one for each clause. A variable group has two vertices for the two values. A clause group has one vertex for each satisfying assignment to its distinct variables. Add edges within groups and between vertices of different groups that disagree on a common variable. If \(m\) is the number of groups, independent sets have size at most \(m\). Size-\(m\) sets choose one vertex per group and correspond bijectively to the models of \(F\). The variable groups ensure this even for unused declared bits.

Write \(v=|V(G)|\), \(t=v+2\) and \(W=2^t-1\). Replace each vertex by \(t\) pairwise nonadjacent twins and each edge by all edges between its two corresponding classes. An independent set using a class has \(W\) possible nonempty intersections with that class. With \(a_i\) independent sets of size \(i\) in \(G\), the resulting graph \(B\) has
\[
 I(B)=\sum_{i=0}^m a_iW^i=AW^m+R,\qquad
 0\le R\le 2^v W^{m-1}<W^m .
\]
The bound uses the total number of subsets of \(V(G)\) and \(m\ge1\). Thus \(I(B)\ge KW^m\) exactly when \(A\ge K\). There are \(vt\) vertices in \(B\) and \(KW^m\) has \(O(r+mt)\) bits.
\end{proof}

\begin{proof}[Completion of Theorem~\ref{main}]
The reductions above establish hardness for both languages. For a formula with \(N\) variables, comparison of its model count with an integer \(K\) in \([1,2^N]\) is in \(\PP\). With one extra leading bit, use the formula on one branch and accept on exactly \(2^N-K+1\) strings in the other. The combined test is strict majority exactly when the count is at least \(K\). For part~1, when \(a\le N\) use \(K=2^{N-a}\). When \(a>N\), positivity of the count suffices, which is 2-SAT. For part~2 use \(K=2^{N-1}\) when \(N\ge1\). A formula with no variables can be evaluated directly.
\end{proof}

\section{An example at equality}

Take \(B\) to be a four-cycle with one diagonal. The empty set, four singletons and the pair on its missing diagonal give \(I(B)=6\). Set \(T=6\), \(b=3\), and take \(H=K_2\), which has \(D=3\) independent sets. The join has \(6+3-1=8\) independent sets. Its edge CNF has six variables, all occurring. Four fresh negative units make \(N=10\) and \(a=7\). Thus \(\Pr[\phi]=8/2^{10}=2^{-7}\), and the selector gives \(\Pr[\psi\wedge L]=1/2\). Adding the missing diagonal makes the source count \(5\) and gives majority probability \(1023/2048\). Removing the remaining diagonal makes the source count \(7\) and gives \(1025/2048\).

\section*{Research disclosure}

\begingroup\hyphenpenalty=10000\exhyphenpenalty=10000\emergencystretch=1em\relax
Under human direction, ChatGPT Astra 6.2 (OpenAI) generated the discovery proof and its research packet. Claude Opus 5.5 (Anthropic) assisted with earlier audits and manuscript preparation. OpenAI Codex performed the mathematical review and revised the paper. GPT-6.2 Sol assisted with literature searches and release preparation. DannyExperiments maintains the repository and is the preferred citation identity. No AI system is proposed as an author.
\par\endgroup

\begin{thebibliography}{9}
\bibitem{akmal}
S. Akmal, \emph{Parameterized Relaxations for Circuits and Graphs}, Ph.D. thesis, MIT, 2024. \href{https://www.shyanakmal.com/publications/ShyanAkmal_PhD-Thesis.pdf}{Full text}.
\bibitem{akmalwilliams}
S. Akmal and R. Williams, \emph{MAJORITY-3SAT (and Related Problems) in Polynomial Time}, 62nd IEEE FOCS, 2021, 1033--1043. \href{https://arxiv.org/abs/2107.02748}{arXiv:2107.02748}.
\bibitem{chanpak}
S. H. Chan and I. Pak, \emph{Computational complexity of counting coincidences}, \href{https://arxiv.org/abs/2308.10214}{arXiv:2308.10214}.
\bibitem{cozman}
F. G. Cozman and D. D. Mau\'a, \emph{The complexity of Bayesian networks specified by propositional and relational languages}, Artif. Intell. \textbf{262} (2018), 96--141. \href{https://arxiv.org/abs/1612.01120}{arXiv:1612.01120}.
\bibitem{finster}
S. Finster, M. Gonz\'alez Amador, E. Lock, F. Marmolejo-Coss\'io, E. Micha and A. D. Procaccia, \emph{Welfare-Maximizing Pooled Testing}, \href{https://arxiv.org/abs/2206.10660v4}{arXiv:2206.10660v4} (2026).
\end{thebibliography}
\end{document}
